\documentclass[sigconf,anonymous,review]{acmart}

%% ---------------------------------------------------------------------------
%% SAC 2027, Semantic Technology track.
%% Double-blind: no names, no affiliation, no repository link in the source.
%% Replace \repourl with an anonymised link before submission, and with the
%% real link in the camera-ready version. It is defined once, here.
%% ---------------------------------------------------------------------------
\newcommand{\repourl}{\textsc{[anonymised repository]}}

\settopmatter{printacmref=false}
\renewcommand\footnotetextcopyrightpermission[1]{}

\usepackage{booktabs}

\begin{document}

\title{Composing Usage Policies for Models Trained on Many Sources}

\begin{abstract}
Machine learning models are increasingly built from data held by several
organisations, each of which attaches its own conditions of use. Standards exist
for expressing those conditions on a single dataset, and a new ODRL profile
standardises actions such as \emph{ai-training} and \emph{ai-augment}. None of
them answer the question that arises as soon as more than one source is
involved: what may be done with the model that results.

We give a composition algebra for ODRL policies and prove it is anti-monotone.
Adding a contributor can never widen what the group may lawfully do with the
model. We ground the algebra in the GA4GH Data Use Ontology, report that one
widely used code cannot be expressed at all because its truth depends on the set
being composed, and measure how far capability falls as a group grows. The rate
at which it falls depends on how similar the contributors' terms are, which
earlier work had to assume; we measure it instead, from the data access
committees of a public archive, and report it as an interval rather than a point.

We then use the algebra to select a group. Because compatibility is not
transitive, this is a search over subsets rather than a filter over candidates,
and anti-monotonicity makes the search tractable: it removes up to 98.6\% of the
space without training anything. In our corpus a group chosen without regard to
policy is permitted to use its own model 20\% of the time. This result follows
from combining declared terms and does not depend on any accuracy measure.
Code and data: \repourl.
\end{abstract}

\begin{CCSXML}
<ccs2012>
<concept>
<concept_id>10002951.10003260.10003282</concept_id>
<concept_desc>Information systems~Web services</concept_desc>
<concept_significance>300</concept_significance>
</concept>
<concept>
<concept_id>10002951.10003227.10003351</concept_id>
<concept_desc>Information systems~Data mining</concept_desc>
<concept_significance>300</concept_significance>
</concept>
<concept>
<concept_id>10010147.10010178</concept_id>
<concept_desc>Computing methodologies~Artificial intelligence</concept_desc>
<concept_significance>300</concept_significance>
</concept>
</ccs2012>
\end{CCSXML}

\ccsdesc[300]{Information systems~Web services}
\ccsdesc[300]{Computing methodologies~Artificial intelligence}

\keywords{ODRL, usage control, data governance, policy composition,
Data Use Ontology, federated learning}

\maketitle

%% ===========================================================================
\section{Introduction}

Large models are more and more often adapted with private data. An institution
fine-tunes a foundation model on its own records, or builds a retrieval system
over a private collection. When several institutions do this together, the
result depends on data from several owners, and each owner attached conditions
to its data.

\subsection{A simple example}
\label{sec:example}

Three people are invited to dinner, and the host must respect every guest's
restriction:

\begin{itemize}
\item guest A eats only fish or vegetables,
\item guest B eats only meat or vegetables,
\item guest C eats only fish or meat.
\end{itemize}

Any two of them can be fed. A and B can share a vegetable dish, A and C can share
fish, B and C can share meat. All three together can eat nothing at all.

The same shape appears in data sharing. Suppose three institutions state the
purposes for which their data may be used:

\begin{itemize}
\item contributor A permits cancer research or method development,
\item contributor B permits neurological research or method development,
\item contributor C permits cancer or neurological research.
\end{itemize}

Any two of them share a purpose. All three share none. A consortium that checked
each member separately against its research question would conclude that every
member is acceptable, and would be wrong.

This small example already shows the three properties that the rest of the paper
makes precise. Conditions \emph{intersect}, so the set of things the group may do
only shrinks as members are added. Adding a member can therefore never add a
permission, which is the anti-monotonicity we prove in
Section~\ref{sec:algebra}. And compatibility is \emph{not transitive}, so
checking members one at a time proves nothing: the problem appears only in the
set of three, which is why Section~\ref{sec:selection} has to search over subsets
rather than filter candidates.

\subsection{The gap}

Standards for expressing conditions exist and are in use.
SPDX~\cite{spdx} describes software licences and supports expressions that
combine them. The GA4GH Data Use Ontology (DUO)~\cite{duo} describes conditions
on biomedical datasets and is published alongside real archives. The W3C Open
Digital Rights Language (ODRL)~\cite{odrl-model,odrl-vocab} gives a general model
of permissions, prohibitions, duties and constraints. In September 2026 the ODRL
community published a draft AI profile~\cite{odrl-ai} that adds actions for AI
workflows: \emph{data-mining}, \emph{ai-training}, \emph{gen-ai},
\emph{ai-augment}, \emph{search} and \emph{index}.

All of these describe what may be done \emph{to} a dataset. None of them say what
happens when several policies must hold at the same time, and none describe the
artefact that training produces. This is the gap we address. The question is not
hypothetical: the new ODRL profile defines \emph{ai-training} for fine-tuning and
\emph{ai-augment} for retrieval and grounding, which are exactly the two ways a
foundation model is adapted with private data today.

\subsection{Scope and contributions}

Choosing a group of contributors has two parts: which groups are
\emph{permitted}, and which of the permitted groups produces the best model. This
paper addresses the first. We assume a utility measure is supplied by the
application, and we make no claim about which permitted group trains best. We
return to this in Section~\ref{sec:selection}, where it turns out that solving the
first part is what makes the second part affordable.

Our contributions are:

\begin{enumerate}
\item A composition algebra for a fragment of ODRL, with a proof that capability
is anti-monotone in the contributor set: adding a contributor can only narrow
what the group may do.
\item A mapping from GA4GH DUO to that fragment, covering all 28 codes observed
in a public archive, with the fidelity of each mapping recorded. One code cannot
be expressed at all, and we explain why.
\item A measurement, rather than an assumption, of how similar contributors'
terms are in practice, taken from the data access committees of the European
Genome-phenome Archive, and reported as an interval.
\item A selection method that uses anti-monotonicity to prune the search over
subsets, and a measurement of what ignoring policy costs: in our corpus a group
chosen without regard to policy may use its own model 20\% of the time.
\item A model bundle that carries its composed agreement, where verification
recomputes the agreement from the recorded declarations instead of comparing a
hash, and where composition is recursive so that reuse can only narrow rights.
\end{enumerate}

%% ===========================================================================
\section{Background and related work}
\label{sec:related}

A good deal of work describes conditions of use, and some of it is widely
deployed. None of it composes conditions from several owners, and none of it
describes the model that training produces. This section sets out what exists,
so that the gap is visible rather than asserted.

\subsection{ODRL and the Data Use Ontology}

ODRL expresses a policy as a set of rules. A \emph{permission} allows an action,
a \emph{prohibition} forbids one, and a \emph{duty} is an obligation that must be
met. Rules carry \emph{constraints}, which restrict when a rule applies, for
example to a purpose, a region or a date.

We use a fragment of ODRL with five actions relevant to building a model:
\emph{use}, \emph{aggregate}, \emph{derive}, \emph{distribute} and
\emph{transfer}. The first three matter most. A model trained on several
contributions is an aggregation of them, and fine-tuning it later is a
derivation. ODRL has both in its core vocabulary, which is why it fits this
problem better than a general access control language. The draft AI
profile~\cite{odrl-ai} adds actions in a hierarchy, where \emph{gen-ai} is a kind
of \emph{ai-training}, which is a kind of \emph{data-mining}. Our fragment can
take those actions directly, because the algebra treats actions through a
subsumption relation rather than as flat labels.

DUO~\cite{duo} is a controlled vocabulary for conditions on biomedical data, used
by the European Genome-phenome Archive (EGA)~\cite{ega} and others. A dataset is
annotated with codes such as HMB (health or medical or biomedical research only),
NCU (no commercial use), or IRB (ethics approval required). We use DUO rather than
invented policies because the distribution of conditions in the wild is itself
part of the result. Inventing a distribution would leave the main measurement
unanswerable: a reader could ask why capability falls the way it does, and the
only honest answer would be that we chose it.

\subsection{Describing conditions on one asset}

DUO has a companion, ADA-M~\cite{adam}, which exists to make data use conditions
automatable and is the closest prior work to ours. Both are built around matching
one request against one dataset. Neither has a notion of combining several
datasets and asking what the joint rights are.

A second group of formats documents datasets and models for human readers.
Datasheets for Datasets~\cite{datasheets} and Model Cards~\cite{modelcards} ask
authors to record provenance, intended use and known limitations.
Croissant~\cite{croissant} gives machine-readable metadata for datasets,
including licence fields, and is used by large dataset repositories. The W3C Data
Privacy Vocabulary~\cite{dpv} supplies controlled terms for purposes, and the ODRL
AI profile uses it for exactly that. These are valuable and they are descriptive:
they record what went in, and define no operation that derives what comes out.

That this matters in practice is not a matter of opinion. An audit of widely used
machine learning datasets found licence information on aggregator sites to be
missing or incorrect at a high rate~\cite{provenance}. Our own measurement points
the same way: of 4\,817 policy records in EGA, only 21.4\% carry machine-readable
conditions at all (Section~\ref{sec:duo}).

\subsection{Licences that travel with derivatives}

Software licence composition is the mature analogue. SPDX~\cite{spdx}, now an ISO
standard, expresses licences and combines them with operators, and copyleft
propagation is the same phenomenon we study: combining incompatible terms yields
something that may not be distributed. Bills of materials turned this into
practice, and CycloneDX~\cite{cyclonedx} extends the idea to models and datasets
with a machine learning bill of materials. These record the ingredients; they do
not derive what may be done with the result.

The OpenRAIL family~\cite{openrail} attaches use restrictions that downstream
users must pass on, and is deployed at scale on widely used models. The mechanism
is contractual and read by people. Nothing computes the result of combining two
such licences with different restriction lists.

C2PA~\cite{c2pa} signs provenance assertions so that the history of a piece of
content can be checked cryptographically. It is the natural prior art for the
signing gap we describe in Section~\ref{sec:bundle}, and it addresses integrity
rather than the meaning of combined terms.

\subsection{Usage control and data spaces}

UCON~\cite{ucon} and the later usage control line study obligations that persist
after access is granted, with decisions that are continuous rather than made
once. That theory underpins our treatment of duties. It is also where ODRL is
most used commercially: industrial data spaces~\cite{ids} apply it to data shared
between companies, with terms such as permitted purpose, retention period and
region. The emphasis there is enforcing terms attached to one asset as it moves
between two parties, rather than combining the terms of many parties over a
shared artefact.

\subsection{Opt-out vocabularies for AI}

A recent group of proposals lets publishers state whether their content may be
used for AI training, including the IETF AI preferences
vocabulary~\cite{aipref} and several industry formats. The draft ODRL AI
profile~\cite{odrl-ai} reconciles them, noting that they are divergent and
fragmented. Regulation is pushing in the same direction: the EU AI
Act~\cite{aiact} requires providers of general-purpose models to summarise
training content, and the Digital Single Market directive~\cite{dsm} allows
rights holders to reserve text and data mining rights in machine-readable form.

All of this is about one publisher and one asset: may you train on this. It is
the input side of the problem.

\subsection{What none of this provides}

The pattern is consistent. Software licence composition is solved, and is about
software. Data use conditions are machine readable, and are matched one dataset
at a time. AI opt-out vocabularies standardise the action, and stop at the asset.
Documentation formats record what went in, and define no operation over it.
Usage control enforces the terms of one agreement, between two parties.

Two things are missing from all of them. There is no semantics for combining the
conditions of several owners, and there is no vocabulary for the artefact that
training produces. Both are needed as soon as a model is built from more than one
governed source, which is now the normal case.

%% ===========================================================================
\section{Composing policies}
\label{sec:algebra}

\subsection{The rules}

Let $S$ be a set of contributors, each with a policy $p_i$. The composed
agreement $A(S)$ is built element by element, as in Table~\ref{tab:rules}.

\begin{table}[h]
\caption{Composition rules. Every rule narrows.}
\label{tab:rules}
\begin{tabular}{ll}
\toprule
Element & Rule \\
\midrule
Permissions & intersect: an action needs every contributor's grant \\
Prohibitions & union: one prohibition binds the whole group \\
Duties & union, merged per action: the strictest applies \\
Constraints & conjoin: set meet, or interval intersection \\
Conflict & deny overrides \\
\bottomrule
\end{tabular}
\end{table}

\subsection{Anti-monotonicity}

Write $\mathrm{Cap}(S)$ for the set of actions the group may perform under
$A(S)$. For contributor sets $S \subseteq S'$,
\[
\mathrm{Cap}(S') \subseteq \mathrm{Cap}(S).
\]

Permissions intersect, so an action permitted by $A(S')$ is permitted by every
member of $S'$ and therefore by every member of $S$. Prohibitions and duties only
accumulate, and constraints only narrow. So admitting a contributor can never
widen what the group may do with the model, however permissive that contributor
is. It can only add data. Removing a contributor can only widen it.

This is the dinner example of Section~\ref{sec:example} stated generally: adding a
guest never adds a dish. We check the property with randomised property-based
testing over generated policies, in addition to the proof, because the
implementation is what the experiments use.

\subsection{Why subsumption is needed}

Constraints must be compared through a subsumption relation and not as strings. A
contributor allowing use in the EU and one allowing use in Germany are
compatible, and their meet is Germany. Under a flat vocabulary they read as a
conflict and the capability is lost.

This is not a detail. If the vocabulary were flat, part of the capability loss we
measure in Section~\ref{sec:cost} would be an artefact of our own modelling
rather than a property of the policies. The same applies to the action hierarchy
of the AI profile, where a contributor permitting \emph{data-mining} and one
permitting only \emph{ai-training} must compose to \emph{ai-training}.

\subsection{Grounding in a real vocabulary}
\label{sec:duo}

We take the distribution of conditions from EGA policy records, which are public
and require no registration. Of 4\,817 policy records, 1\,031 carry structured
DUO codes: an adoption rate of \textbf{21.4\%}. That number is a result about
practice rather than an inconvenience. Most data use conditions in this archive
are free text or an external link, and are not machine readable at all.

We classify all 28 observed codes by how they map into the fragment
(Table~\ref{tab:duo}). Three are deprecated: retired from the ontology but still
present in live annotations, so a tool that meets real governance metadata has to
decide what to do with obsolete terms.

\begin{table}[h]
\caption{Classification of the 28 DUO codes observed in EGA.}
\label{tab:duo}
\begin{tabular}{llr}
\toprule
Class & Meaning & $n$ \\
\midrule
permission & sets the purpose the data may serve & 8 \\
restriction & a constraint or prohibition & 8 \\
externally adjudicated & an authority decides out of band & 4 \\
deprecated & obsolete, still present in annotations & 3 \\
duty & an obligation rather than a restriction & 2 \\
structural & an ontology category, not a condition & 2 \\
out of fragment & cannot be expressed & 1 \\
\bottomrule
\end{tabular}
\end{table}

One code cannot be expressed at all. DUO:0000020 (COL, collaboration required)
means the data provider must be a participant in the work. Whether it is
satisfied therefore depends on the very contributor set being composed. That is
self-referential, and it is the one construct our fragment excludes by design. We
report it rather than hide it: a widely used standard contains a condition that a
compositional treatment cannot accept without a fixed point, and anyone building
on DUO will meet the same problem.

Every mapping carries a fidelity marker, \emph{exact} or \emph{interpretive}, so
the places where DUO and ODRL do not line up cleanly stay visible to a reader.

%% ===========================================================================
\section{What composition costs}
\label{sec:cost}

We draw contributor policies from the observed distribution of DUO code sets and
compose subsets of increasing size, measuring the share of groups that retain any
usable permission.

The result depends on one modelling choice: how similar the contributors' terms
are. A group of similar institutions loses little; a group of unrelated ones
loses a great deal. EGA gives real code frequencies and real co-occurrence, but
earlier versions of this study had to \emph{assume} which institution holds which
conditions, and the assumption carried the result. Across seven assumed models
the loss at ten contributors ranged from 7\% to 77\%, which is most of the
possible range.

\subsection{Measuring similarity instead of assuming it}

EGA policy records name the Data Access Committee that issues them. Grouping
policies by committee makes similarity observable.

\paragraph{What the similarity number means.} We describe similarity with a
single number $c$, between 0 and 1. It answers one question: if two institutions
belong to the same group, how likely are they to have written the same
conditions? At $c=0$ every institution decides on its own, and two of them match
only by accident. At $c=1$ every institution follows the same template, and any
two of them always match. Real groups lie between.

This number is what decides how quickly capability falls. If all contributors
have identical conditions, combining them costs nothing, because intersecting a
set with itself returns the same set. If their conditions differ, every
additional contributor removes something. So $c$, and not the number of
contributors, is what governs the result.

\paragraph{Why it cannot simply be counted.} We cannot observe $c$ directly. We
can only observe how often two policies happen to be identical, and that is not
the same quantity, because two institutions can agree by coincidence: some
condition sets are common, so two institutions deciding independently will
sometimes write the same thing. The observed rate therefore mixes two causes,

\[
\text{policies match} \;=\;
\underbrace{\text{both followed a shared template}}_{\text{similarity}}
\;+\;
\underbrace{\text{both decided alone and agreed}}_{\text{coincidence}} .
\]

The model below separates them. Write $p$ for the coincidence rate, that is, how
often two independent draws from the observed distribution of code sets produce
the same set. For the EGA distribution $p = 0.0131$, so two unrelated
institutions agree about 1.3\% of the time for no reason other than chance. Any
agreement above that rate is evidence of shared practice, and the model converts
the excess into a value of $c$.

\paragraph{The model.} For two contributors, writing $p$ for the chance that two
independent draws coincide, the probability that their code sets are identical is
\[
\mathrm{agreement} = c^2 + (1-c)^2 p + 2c(1-c)w,
\]
where $w$ is the weight of the shared set. The shared set is itself drawn from
the same distribution, so $\mathbb{E}[w]=p$ and this reduces to
\[
\mathrm{agreement} = p + c^2 (1-p),
\qquad
c = \sqrt{\frac{\mathrm{agreement}-p}{1-p}}.
\]

Measuring agreement in EGA then gives $c$, as shown in
Table~\ref{tab:cohesion}. The two rows read very differently. Policies issued by
two \emph{different} committees are identical 1.18\% of the time, which is
slightly below the 1.31\% that chance alone would produce, so there is no
evidence of shared practice between institutions and $c$ is zero. Policies issued
by the \emph{same} committee are identical 47.4\% of the time, far above chance,
giving $c = 0.683$.

\paragraph{Why an interval and not a single number.} A committee is not a
consortium. It is a single institution issuing several policies, usually written
by the same people under the same internal rules, so it is the most similar
grouping that can be found anywhere in the archive. A real consortium is made of
\emph{different} institutions and must therefore be less similar than that. The
same-committee figure is consequently an upper bound on $c$, and the
across-committee figure a lower bound. The truth for any particular consortium
lies somewhere between, and we report the whole interval rather than choose a
point inside it, because choosing would mean replacing the assumption we set out
to remove with a different one.

\begin{table}[h]
\caption{Similarity measured from EGA, over 1\,031 annotated policies and 899
committees.}
\label{tab:cohesion}
\begin{tabular}{lrl}
\toprule
Quantity & Value & Implies \\
\midrule
chance coincidence $p$ & 0.0131 & --- \\
agreement, different committees & 0.0118 & $c=0.000$ \\
agreement, same committee & 0.474 & $c=0.683$ \\
\bottomrule
\end{tabular}
\end{table}

Policies from different committees agree slightly \emph{less} often than chance,
which this model cannot express at all, since it can only add similarity. The
lower bound is therefore zero.

\subsection{How much capability is lost}

Table~\ref{tab:erosion} re-runs the study at both ends of the measured interval,
and with a corpus that assumes nothing: code sets drawn whole from the committees
that published them, keeping whatever structure EGA actually has.

\begin{table}[h]
\caption{Capability lost at ten contributors. The assumed models span 7\%--77\%.}
\label{tab:erosion}
\begin{tabular}{lrr}
\toprule
Condition & Loss at 10 & Retains any \\
\midrule
measured ceiling ($c=0.683$) & 16\% & 84\% \\
\textbf{drawn from real committees} & \textbf{52\%} & \textbf{48\%} \\
measured floor ($c=0.000$) & 53\% & 47\% \\
\midrule
assumed baseline, for reference & 64\% & 36\% \\
\bottomrule
\end{tabular}
\end{table}

Measuring the assignment narrows the range from 69 points wide to 37. The loss is
large under every condition, so the direction of the result does not depend on
the modelling choice, but the magnitude does and must always be quoted with the
condition that produced it.

The upper bound rests on the 57 committees that publish more than one
machine-readable policy. That is a small number, and it is also a finding:
alongside the 21.4\% adoption rate, it says that structured data use policy is not
yet common practice.

\subsection{Selecting a group}
\label{sec:selection}

If composing an arbitrary group usually yields a model nobody may use, the
practical question is how to find a group that does.

As Section~\ref{sec:example} showed, this is a search rather than a filter.
Compatibility is not transitive, so feasibility is a property of the set, and
candidates cannot be checked one at a time against a reference policy.

Anti-monotonicity makes the search tractable. Because capability only narrows as
the set grows, feasibility is downward closed: every subset of a feasible set is
feasible. So a contributor that fails alone can be dropped outright, and a set of
size $k+1$ need only be considered when all its subsets of size $k$ are feasible.
This is exactly the candidate generation of the Apriori algorithm~\cite{apriori},
and the pruning is sound because the property it rests on is exact rather than
estimated.

\paragraph{The baseline.} The comparison needs a group chosen the way one would
be chosen today. Our policy-blind baseline is therefore not a random group:
contributors are ranked by how much data they bring, and the best $N$ are taken,
with $N$ set to the size the policy-aware search returned. Both groups then have
the same number of members and differ only in whether policy was consulted. This
is deliberately the strongest version of the baseline, because ranking
contributors by the data they hold is standard practice, and a weaker baseline
would make the comparison easier than it should be.

\begin{table}[h]
\caption{Selection over a pool of 16, averaged over 5 corpora. The policy-blind
group is the best $N$ contributors by data volume, with $N$ matched to the size
of the policy-aware group. ``Blind usable'' is how often that group is permitted
to use the model it produced.}
\label{tab:selection}
\begin{tabular}{lrr}
\toprule
Request & Space removed & Blind usable \\
\midrule
actions only & 60.9\% & 20\% \\
purpose = cancer research & 83.7\% & 0\% \\
purpose = research & 98.6\% & 0\% \\
\bottomrule
\end{tabular}
\end{table}

Two things in Table~\ref{tab:selection} are worth reading together. The pruning
works best exactly where selection is hard, and degenerates towards full
enumeration where the requirement is so loose that almost any group would have
satisfied it. That is a property of the problem and not a weakness of the
algorithm: when the constraint is weak, the family of feasible sets really is that
large.

And a group assembled without regard to policy is usually not permitted to use the
model it produced. Under the loosest request it is permitted 20\% of the time;
under a request naming a purpose, never. A policy-aware group is permitted by
construction. This result follows from combining declared terms. It never touches
a model, a metric or a dataset, so no accuracy measurement can change it, which is
why we report it as settled while making no claim about model quality.

\subsection{What the legal filter makes affordable}

The filter is exact and costs nothing: no training, and no data leaves any site.
So it can run first, and whatever method is used afterwards to rank the remaining
groups by expected model quality only has to consider groups that could actually
be deployed. Training a group that may not use its own model is wasted
computation. The order is not a convenience; it is the only order that makes
sense.

We do not rank the permitted groups. That is not the same as claiming they are
equally good, and the distinction matters: it means only that choosing among them
needs a measure of expected model quality, which this paper does not provide and
does not assume. What the filter guarantees is that whatever measure is used
later will be applied to groups that could lawfully be deployed. We return to
this in Section~\ref{sec:future}.

%% ===========================================================================
\section{Models that carry their policy}
\label{sec:bundle}

A trained model is a derivative of every contributor's data, and the composed
agreement states what may lawfully be done with it. Separate the two and the
terms are lost the first time the file is copied. We export a bundle holding the
weights and a manifest: the agreement, every contributor declaration it was
composed from, and what remains owed.

\subsection{Verification recomputes the agreement}

Checking a bundle does not compare a digest against itself. It recomposes the
agreement from the declarations recorded in the manifest and confirms that the
result matches what the manifest claims, then checks that the weights still hash
to what was recorded. Widening a purpose therefore requires forging every
contributor's declaration, not editing one line.

\subsection{Composition is recursive}

An agreement is itself a policy, so a bundle's agreement composes back in when the
model is reused as an input. Capability can only narrow, which means a restricted
model cannot be turned into an unrestricted one by training something else on it.
This is the property that makes a policy-carrying model more than documentation,
and it is the step that none of the formats in Section~\ref{sec:related} take.

\subsection{What the bundle is not}

The manifest is not signed. A forger who rewrites the declarations and the
agreement consistently produces a bundle that verifies. Detection covers accident
and one-sided edits; closing the rest needs signatures over the digest and a key
distribution story, for which C2PA~\cite{c2pa} is the obvious model. We state this
rather than leave it to be discovered.

%% ===========================================================================
\section{Discussion}

\subsection{Limitations}

Bundles are not signed, as just described.

The similarity upper bound rests on 57 committees. It is reported with a
bootstrap interval, and the lower bound is a hard floor rather than an estimate.

Model quality is not addressed. We make no claim about which permitted group
produces the better model. Our evaluation of capability never touches a metric,
which is what makes it stable, but it also means this paper says nothing about
accuracy.

We assume contributors' data is of comparable quality. The legal axis cannot see
data quality, so a contributor with permissive terms and poor data is invisible to
the method described here and would have to be excluded by some other means.

Most of this is not specific to federated learning. The algebra applies to any
derivative built from several governed sources: a trusted research environment, a
shared analysis platform, or a model fine-tuned on pooled records. Federated
learning~\cite{fedavg} is where we implemented it, because there the policy check
can run inside the system before training starts, not because the result depends
on it. Centralising the data does not avoid the problem: the model is still a
derivative of every source, and every source's terms still bind it.

\subsection{Ranking the permitted groups}
\label{sec:future}

This is the obvious next step and the one we did not take. Several directions
exist. Contribution measures such as Data Shapley~\cite{datashapley}, influence
functions or leave-one-out estimate what each contributor adds, at a cost that
grows quickly with group size. Client selection methods from federated learning
use cheaper signals: Oort~\cite{oort} ranks participants by a statistical utility
derived from their local loss. Short exploratory runs are a third option: train
every permitted group for a few rounds, rank them, and continue only with the
best.

One asymmetry makes this harder than it looks, and it is worth stating because it
shapes what is possible. Permissions can be taken from a declaration, because
claiming more restrictions than you have only shrinks what you may do: there is no
incentive to overstate them. Data quality cannot, because the incentive runs the
other way. Quality has to be measured, and measuring it is the expensive problem
that the legal filter makes affordable rather than solves.

\subsection{Conclusion}

Conditions on data compose, and they compose in one direction. We gave an algebra
for a fragment of ODRL, proved that capability is anti-monotone in the contributor
set, grounded it in a real vocabulary, and measured what it costs: a group
assembled without regard to policy is usually not permitted to use the model it
produced. We also measured the one quantity earlier versions of this study had to
assume, and report it as an interval rather than a point.

Signing the manifest is the remaining open item, and the one that would turn
policy-carrying models from a description into a mechanism.

\subsection*{Availability}

The implementation, the policy corpus, the committee metadata and every result
table in this paper are available at \repourl. All tables are produced from
committed result files, so each number here can be traced to the run that produced
it.

\bibliographystyle{ACM-Reference-Format}
\bibliography{refs}

\end{document}
